\scanpage{51}
\margin{50}
\mspage{5}

\deleted{
\[
(\unsure{20|19})\qquad\uline{\bar\omega}\xrightarrow{\sim}\{\pm1\}\quad\text{\unsure{i.e.} on dispose }\uline{\omega}\xrightarrow{\sim}\uline{\varepsilon}
\]
donc dorénavant on identifie l'ens $\uline{\omega}=\Omega(J)$ \inserted{$=\Omega(\Sigma_{\mathbb R})$} avec l'ens $\uline{\varepsilon}$ des deux hémisphères.
}

\margin{Il s'agit d'\uline{arcs} de \unsure{l'un des} cercles, \unsure{qui sont déterminés} \unsure{par} origine, extrémité, désignation du parcours à priori \unsure{court|et} \unsure{ou|en} un des \unsure{sens}}

\margin{\uline{NB}
\[
\operatorname{Rep}(J)\simeq J\times\uline{\omega}
\]
\[
(i,j,k)\longmapsto(i,\omega_{i\to j\to k})\qquad(\unsure{20})
\]
le paquet $(a_i^{\omega\varepsilon})$ est un torseur sous $\Gamma=G_J\times\{1,\tau\}$
}

Introduisons sur $\widetilde\Sigma$ les \deleted{\unsure{arcs}} chemins simples : (segments orientés plongés)

\[
(\unsure{21})\qquad\overset{12}{\overset{\rule[-.2ex]{.4pt}{.6ex}\rule[.4ex]{2em}{.4pt}\rule[-.2ex]{.4pt}{.6ex}}{a_i^{\omega,\varepsilon}}},\quad\overset{6}{\overset{\rule[-.2ex]{.4pt}{.6ex}\rule[.4ex]{1.3em}{.4pt}\rule[-.2ex]{.4pt}{.6ex}}{b_i^\omega}},\quad\overset{6}{\overset{\rule[-.2ex]{.4pt}{.6ex}\rule[.4ex]{1.3em}{.4pt}\rule[-.2ex]{.4pt}{.6ex}}{c_i^\varepsilon}},\quad\overset{6}{\overset{\rule[-.2ex]{.4pt}{.6ex}\rule[.4ex]{1.3em}{.4pt}\rule[-.2ex]{.4pt}{.6ex}}{d_i^\varepsilon}},\qquad
i\in J,\ \omega\in\uline{\omega},\ \varepsilon\in\uline{\varepsilon}
\]

chaque paquet \inserted{$b,c,d$} formant un torseur sous $G_J$

\margin{\uline{NB} on a donc \deleted{\unsure{déduit}}
$g(a_i^\omega)=a_{g(i)}^{g(\omega)}$ etc}

$\tau$ opérant par
\[
(22)\quad \tau a_i^{\omega,\varepsilon}=a_i^{\omega,\unsure{\varepsilon'}},\quad
\tau b_i^\omega=b_i^\omega,\quad \tau c_i^\varepsilon=c_i^{\varepsilon'},\quad
\tau d_i^\varepsilon=d_i^{\varepsilon'}\;\raisebox{.4ex}{$\scriptstyle\unsure{\nearrow}$}.
\]

On a
\[
(23)\qquad a_i^{\omega\varepsilon}:{\unsure{S}}_i^\varepsilon\longrightarrow\deleted{\unsure{S_i^\varepsilon}}R_i^\omega\text{ sur }U_i\quad\inserted{(\text{quart de tour})}
\]
\deleted{dans le sens de l'orientation de $\Sigma$ en $R_i$}

ce qui caractérise $a_i^{\omega\varepsilon}$. \deleted{\unsure{explicite le fait q}}

\[
(24)\qquad b_i^\omega:Q_i\longrightarrow R_{\omega'i}^{\omega'}\text{ sur }\widetilde\Sigma_{\mathbb R}\simeq\deleted{\unsure{\Sigma_{\mathbb R}}}\inserted{\operatorname{\text{Déc}}(\Sigma_{\mathbb R},J)}\ \deleted{\unsure{\text{sur les}}}
\]
\deleted{(chemin d'origine \unsure{sur ce cercle}}

(chemin sur $\widetilde\Sigma_{\mathbb R}$, d'origine $Q_i$, dans la direction $\omega$, jusqu'à l'extrémité \deleted{\unsure{qui est}} du segment composante de $\widetilde\Sigma_{\mathbb R}\simeq\operatorname{\text{\unsure{Déc|Dic}}}(\Sigma_{\mathbb R},J)$, savoir $R_{\omega'i}^{\omega'}$.

\[
(25)\qquad c_i^\omega:P^\omega\mathrel{-\!\!\!-\!\!\!-}Q_i\qquad\text{suivant un arc de grand cercle}
\]

\[
(26)\qquad d_i^\omega:P^\omega\longrightarrow S_i^\omega\qquad\text{suivant un arc de grand cercle}
\]

\margin{
\[
\left.\begin{array}{l}\hline
\text{On pose aussi}\\
(27)\quad\widetilde b_i^\omega:R_{\omega i}^{\omega}\longrightarrow R_{\omega'i}^{\omega'},\quad\widetilde b_i^\omega=b_i^\omega(b_i^{\omega'})^{-1}\\\hline
\end{array}\right)
\]
\[
\left.\begin{array}{l}
\text{chemins dans}\\
\text{la direction }\omega)\\
\text{\unsure{et non pour}}\ \unsure{b_i}\\
\text{\deleted{\unsure{des points}}}\\
\text{\deleted{\illegible{2}}}\ \deleted{\unsure{b_i^\omega}}
\end{array}\right)
\]
}

Tous les arcs $\unsure{\widetilde b_i^\omega|\bar b_i^\omega},c_i^\omega,d_i^\omega$ se font sur des grands cercles de $\Sigma$, sauf que les extrémités, des $c_i^\omega,d_i^\omega$ ne sont pas dans $\Sigma^*$ mais dans \inserted{$U_i\subset$} $\partial\widetilde\Sigma\setminus\Sigma^*$
